\section{An integer in the conductor}
\label{sec:conductor}

Let $T$ be a unital subring of $D$ in which some nonzero element admits arbitrarily long decompositions into nonunits.
By Proposition~\ref{prop:generic}, $\mathbb{Q}T = B$.
In particular, some positive integer $n$ satisfies $nt \in T$.
We will prove that every such $n$ satisfies $nD \subseteq T$.

Fix a rational prime $p$ and put
\[
    T' = \mathbb{Z}_p \otimes_{\mathbb{Z}} T,
    \qquad
    D' = \mathbb{Z}_p \otimes_{\mathbb{Z}} D,
    \qquad
    B' = \mathbb{Q}_p \otimes_{\mathbb{Q}} B.
\]
The ring $B'$ is the localization of $\mathbb{Q}_p[t]$ obtained by inverting the $F_i$.
These are ordinary tensor products, so their elements are finite sums of pure tensors.
The module $\mathbb{Z}_p$ is flat over $\mathbb{Z}$ because it is torsion-free over a PID.
Tensoring $T \subseteq D \subseteq B$ therefore gives inclusions, and $\mathbb{Z}_p \otimes_{\mathbb{Z}} B = B'$ since $B$ is a $\mathbb{Q}$-algebra.
Moreover $T'[1/p] = \mathbb{Q}_p \otimes_{\mathbb{Q}} (\mathbb{Q}T) = B'$.
Neither $\alpha_p$ nor $\beta_p$ is a root of any $F_i$, since both arguments are transcendental over $\mathbb{Q}$.
Consequently evaluation at either argument extends to $B'$.
At $\gamma \in \{\alpha_p, \beta_p\}$ the value on $\sum_j u_j \otimes f_j \in D'$ is $\sum_j u_j f_j(\gamma) \in \mathbb{Z}_p$.
The scalar extension puts $\gamma$ in the coefficient field, so evaluation on $B'$ now has kernel $(t - \gamma)B'$.
The kernel of the pair of evaluations is consequently generated by $L_p(t) = (t - \alpha_p)(t - \beta_p)$.
If this entire kernel lies in $T'$, membership in $T'$ can be checked on the two values alone.

By Lemma~\ref{lem:extraction}, there are words in $T$ from arbitrarily late stages $i$ of the construction, of the form
\[
    w_i = \frac{F_i(t)^{e_i}}{H_i(t)},
    \qquad 1 \le e_i \le R,
\]
where $R = 2r$ is fixed and a possible sign has been absorbed using $-1 \in T$.
Here $r$ bounds the primary valuations of the element with unbounded decomposition lengths, and $H_i$ denotes the denominator of the particular word selected at stage $i$.

We will use the following test twice.
For $g \in B'$, the equality $B' = T'[1/p]$ gives
\[
    g B' \subseteq T'
    \quad\Longleftrightarrow\quad
    g/p^k \in T' \quad\text{for every } k \ge 0.
\]
Indeed every element of $B'$ is $u/p^k$ for some $u \in T'$.
We first apply the test to a power of $L_p$.
We then show that any repeated factor of the conductor generator could be cancelled, by taking a derivative at its root.

\begin{lemma}[Exact polynomial transfer]
    \label{lem:exact-transfer}
    One has $L_p(t)^R B' \subseteq T'$.
\end{lemma}

\begin{proof}
    Fix a positive integer $n$ with $nt \in T$ and put $s = \nu_p(n)$.
    The scalar $n/p^s$ is a unit in $\mathbb{Z}_p$, so $p^s t \in T'$ and $A_s = \mathbb{Z}_p[p^s t] \subseteq T'$.
    Put $g = p^{2s}L_p(t) \in A_s$.
    Fix any integer $k \ge 0$ and choose an extracted stage $i \ge k$ sufficiently late for \eqref{eq:transfer-module}.
    As $w_i \in T'$, that estimate gives $g^{e_i}/p^i \in A_s + A_s w_i \subseteq T'$.
    Multiplication by $p^{i - k}g^{R - e_i} \in T'$ yields $g^R/p^k \in T'$.
    The test above puts $g^R B'$ in $T'$, and $g^R B' = L_p^R B'$ since $p$ is invertible in $B'$.
\end{proof}

\begin{lemma}[The local conductor]
    \label{lem:local-conductor}
    The conductor $(T' : B') = \{f \in B' \mid fB' \subseteq T'\}$ equals $L_p(t)B'$.
\end{lemma}

\begin{proof}
    Write $I = (T' : B')$, the largest ideal of $B'$ contained in $T'$.
    Lemma~\ref{lem:exact-transfer} gives $L_p^R B' \subseteq I$.
    For $f \in I$, every $p^{-k} f$ with $k \ge 0$ belongs to $T'$ and has integral values at $\alpha_p$ and $\beta_p$.
    Both values of $f$ are therefore zero.
    The two distinct linear factors of $L_p$ survive in the localized principal ideal domain $B'$, so
    \[
        L_p^R B' \subseteq I \subseteq L_p B'.
    \]
    Thus we may write $I = h B'$ where $h$ is a product of powers of $t - \alpha_p$ and $t - \beta_p$, each with exponent between one and $R$.

    Suppose $(t - \gamma)^2$ divides $h$ at one of the two arguments.
    Put $q = h/(t - \gamma)^2$.
    We will prove $(h/(t - \gamma)) B' \subseteq T'$, contradicting the maximality of $I$.
    Choose an integer $d \ge 0$ with $p^d q \in T'$, using $T'[1/p] = B'$.
    Fix $k \ge 0$.
    Along the extracted sequence, $\nu_p(w_i'(\gamma)) \to -\infty$ by Proposition~\ref{prop:construction}.
    We may therefore choose $w_i$ with $u = 1/(p^k w_i'(\gamma)) \in p^d\mathbb{Z}_p$.
    Then $uq \in T'$, and the scalar $w_i(\gamma)$ belongs to $\mathbb{Z}_p \subseteq T'$.

    In $B'$ the first-order expansion at $\gamma$ gives
    \[
        w_i(t) - w_i(\gamma) \equiv w_i'(\gamma)(t - \gamma)
        \pmod{(t - \gamma)^2 B'}.
    \]
    To see this for a rational function, expand its numerator and denominator modulo $(t - \gamma)^2$ and invert the denominator, whose value at $\gamma$ is nonzero.
    After multiplication by $uq$ the error belongs to $h B' = I$.
    Consequently
    \[
        uq\bigl(w_i(t) - w_i(\gamma)\bigr)
        \equiv \frac{h/(t - \gamma)}{p^k} \pmod{I}.
    \]
    The left side belongs to $T'$ and $I \subseteq T'$, so $(h/(t - \gamma))/p^k \in T'$ for every $k$.
    The test above now puts $(h/(t - \gamma))B'$ in $T'$.
    This ideal strictly contains $I$ because $t - \gamma$ is a nonunit of $B'$.
    Neither factor can be repeated, so $I = L_p B'$.
\end{proof}

\begin{theorem}[The integer conductor]
    \label{thm:conductor}
    Suppose a unital subring $T \subseteq D$ has a nonzero element with unbounded decomposition lengths.
    There exists a positive integer $n$ with $nt \in T$, and every such $n$ satisfies $nD \subseteq T$.
\end{theorem}

\begin{proof}
    Existence follows from $\mathbb{Q}T = B$.
    Fix any such $n$ and any rational prime $p$.
    By Lemma~\ref{lem:local-conductor} the kernel of the value map $B' \to \mathbb{Q}_p^2$ is contained in $T'$.
    The value vectors of $1$ and $t$ form the columns of the matrix
    \[
        \begin{pmatrix}
            1 & \alpha_p \\
            1 & \beta_p
        \end{pmatrix}.
    \]
    Its determinant is a $p$-adic unit because $\alpha_p \in p\mathbb{Z}_p$ and $\beta_p \in \mathbb{Z}_p^\times$.
    Since $n, nt \in T'$, the image of $T'$ contains $n\mathbb{Z}_p^2$.
    Each element of $nD'$ has its two values in $n\mathbb{Z}_p^2$ and hence has the same values as some element of $T'$.
    Their difference lies in the kernel, so $nD' \subseteq T'$.

    For descent, take the abelian group $M = (nD + T)/T$.
    Flatness of $\mathbb{Z}_p$ over $\mathbb{Z}$ identifies $\mathbb{Z}_p \otimes_{\mathbb{Z}} M$ with $(nD' + T')/T' = 0$.
    To see that the tests at all primes detect $M$, observe that $M$ is torsion since $\mathbb{Q}T = B$.
    A nonzero element of $M$ would generate a subgroup isomorphic to $\mathbb{Z}/m\mathbb{Z}$ with $m > 1$.
    For a prime $p \mid m$, flatness would embed the nonzero group $\mathbb{Z}_p/m\mathbb{Z}_p$ in $\mathbb{Z}_p \otimes_{\mathbb{Z}} M$.
    Hence $M = 0$, which is the asserted inclusion.
\end{proof}
